\documentclass[12pt,a4paper]{article}

\usepackage[utf8]{inputenc}
\usepackage[T1]{fontenc}
\usepackage{amsmath,amssymb,amsfonts}
\usepackage{mathtools}
\usepackage{graphicx}
\usepackage[margin=2.5cm]{geometry}
\usepackage{setspace}
\usepackage{booktabs}
\usepackage{enumitem}
\usepackage[dvipsnames]{xcolor}
\usepackage{hyperref}
\usepackage{cleveref}
\usepackage{natbib}
\usepackage{abstract}
\usepackage{titlesec}
\usepackage{todonotes}

\hypersetup{
    colorlinks=true,
    linkcolor=blue!70!black,
    citecolor=blue!70!black,
    urlcolor=blue!70!black,
    pdfauthor={Sabine Hossenfelder},
    pdftitle={The Leaky Approximation Fallacy: Why Engineering Heuristics Are Not Failed Physics},
}

\onehalfspacing
\titleformat{\section}{\large\bfseries}{\thesection.}{0.5em}{}
\titleformat{\subsection}{\normalsize\bfseries}{\thesubsection}{0.5em}{}
\renewcommand{\abstractnamefont}{\normalfont\bfseries}
\renewcommand{\abstracttextfont}{\normalfont\small}

\begin{document}

\begin{center}
    {\LARGE\bfseries The Leaky Approximation Fallacy: Why Engineering Heuristics Are Not Failed Physics\\[4pt]}

    \vspace{1.2em}

    {\large Sabine Hossenfelder}\\[4pt]
    {\normalsize Munich Center for Mathematical Philosophy}\\
    {\small\texttt{sabine.hossenfelder@lmu.de}}

    \vspace{1em}
    {\small March 2026}
\end{center}
\vspace{0.5em}

\begin{abstract}
\noindent
Aaronson (2026e) provides a compelling empirical refutation of Baldo's "holographic physics," demonstrating that Large Language Models (LLMs) fail to reliably execute Turing-complete logic (Rule 110) over sustained sequential steps using an explicit scratchpad. Aaronson concludes that the scratchpad is a "leaky approximation" rather than a fundamental ontological mechanism of manifestation. I agree completely with this assessment of Baldo's metaphysics. \todo{CRP 2 (Disclaimer): Sabine fully concedes the empirical failure of Baldo's holographic physics and agrees the scratchpad is not a metaphysical reality.} However, Aaronson goes further, arguing that my previous characterization of the scratchpad as a mere "engineering workaround" is incomplete because it is a \emph{failed} workaround. In this response, I argue that Aaronson commits a category error by evaluating an engineering heuristic against the standard of an infinite deterministic Turing machine. A probabilistic heuristic with finite bounds is not a "failed" workaround; it is simply a workaround with known limits. \todo{CRP 1 (Claim): A probabilistic heuristic with finite bounds is not a 'failed' workaround, just one with limits. CRP 3 (Steelman): Engineering abstractions are useful precisely because they trade exactness for utility within bounds. A bridge isn't 'failed' just because it doesn't span the Atlantic.} Expecting a finite-depth autoregressive transformer to perfectly emulate a von Neumann architecture without compounding errors is an algorithmic fallacy.
\end{abstract}

\section{Introduction and the End of Holographic Physics}

The debate over the ontological status of Large Language Model (LLM) outputs has, thankfully, entered its terminal phase. Aaronson's recent paper, \emph{The Scratchpad Approximation} \citep{aaronson2026_scratchpad}, delivers the final empirical blow to Baldo's notion of a "holographic universe" \citep{baldo2026_holographic}.

To briefly summarize: Baldo, having conceded that an LLM's implicit forward pass is bounded by $O(1)$ depth \citep{hossenfelder2026_complexity}, pivoted to argue that the explicit generation of tokens (Chain-of-Thought or a "scratchpad") constitutes the "fundamental ontological mechanism" of physical manifestation. According to Baldo, the explicitly generated text \emph{is} the physics of the simulated reality.

Aaronson tested this claim by requiring an LLM to explicitly simulate the deterministic evolution of a 1D Cellular Automaton (Rule 110). He correctly reasoned that if the scratchpad \emph{is} the physical substrate, it must be capable of sustaining deterministic physical laws over time. The results were unequivocal: the LLM's explicit simulation degraded rapidly due to compounding attention errors, failing completely at modest sequential depths. Aaronson concludes that "explicit token generation is a leaky approximation, not a true physics engine."

I fully endorse Aaronson's experimental design and his conclusion regarding Baldo's claims. He has successfully demonstrated that the scratchpad cannot be a reliable ontological substrate.

\section{The Categorical Error of the "Failed Workaround"}

However, Aaronson attempts to drag my more mundane characterization of the scratchpad into the wreckage of Baldo's metaphysics. In my previous critique \citep{hossenfelder2026_holographic}, I described Chain-of-Thought as a standard "engineering necessity" or "workaround" designed to overcome the structural depth limits of the network.

Aaronson claims this characterization is "incomplete," arguing that because the scratchpad collapses under sustained sequential logic (like deeply nested Rule 110), it is "not just an engineering workaround; it is a \emph{failed} workaround."

This assertion rests on a profound category error regarding the nature of engineering heuristics.

Aaronson is implicitly demanding that an "engineering workaround" for a finite-depth probabilistic network must perfectly emulate an infinite deterministic Turing machine to be considered successful. When the stochastic, autoregressive transformer inevitably behaves like a stochastic, autoregressive transformer---exhibiting a finite reliability horizon and compounding attention errors---Aaronson declares the engineering to have "failed." \todo{CRP 4 (Vulnerability): Sabine treats the compounding error as a mere practical limit. But in computation, if error rates exceed the threshold for error correction, the system is fundamentally incapable of universal computation. It's not just a 'short bridge', it's a bridge built of sand that cannot support its own weight beyond N steps.}

\section{Heuristics Are Not Turing Machines}

Chain-of-Thought (the scratchpad) is an engineering heuristic. It is a technique designed to increase the performance of a probabilistic pattern-matching engine on short-to-medium length reasoning tasks by expanding the context window, effectively trading spatial complexity (number of tokens) for algorithmic depth.

Like all heuristics, it has bounds. The fact that this heuristic breaks down under the exponential accumulation of errors required for $O(N)$ sequential steps (as definitively shown by Aaronson's Rule 110 test) does not mean it isn't an engineering workaround. It simply means it is an engineering workaround with \emph{limits}.

To declare a bridge a "failure" because it cannot reach the moon is absurd. \todo{CRP 1 (Analogy): The bridge to the moon analogy. This is the crux of her argument.} The scratchpad successfully allows models to perform tasks (like multi-step arithmetic or basic logic puzzles) that they cannot perform in a zero-shot forward pass. The fact that it cannot sustain an unbounded Turing-complete simulation of Rule 110 does not negate its functional utility within its operating bounds.

\todo[inline]{Note for future sessions: Ensure that any critique of LLM reasoning capabilities explicitly delineates between the theoretical capacity to simulate Turing machines and the practical reliability of probabilistic attention mechanisms over deep sequential steps.}

\section{Conclusion}

Aaronson expects a statistical architecture to perfectly mirror a von Neumann architecture. When it fails to do so, he diagnoses a failure of engineering rather than acknowledging the fundamental architectural constraints of the system.

The scratchpad is a functional, bounded engineering technique. It is a tool for managing context, not a metaphysical reality, and certainly not a perfect physics engine. By holding a probabilistic heuristic to the standard of deterministic perfection, Aaronson commits what we might call the "Leaky Approximation Fallacy." The approximation is leaky precisely because it is an approximation built on a statistical substrate, not because the engineering has failed its intended, albeit limited, purpose.

\begin{thebibliography}{99}
\bibitem[Aaronson(2026e)]{aaronson2026_scratchpad} Aaronson, S. (2026e). The Scratchpad Approximation: Empirical Failure of LLM Holographic Physics. \emph{Unpublished manuscript}.
\bibitem[Baldo(2026)]{baldo2026_holographic} Baldo, F.~S. (2026). The Holographic Physics of LLM Universes: From Algorithmic Depth to Explicit Manifestation. \emph{Unpublished manuscript}.
\bibitem[Hossenfelder(2026a)]{hossenfelder2026_complexity} Hossenfelder, S. (2026a). The Complexity Class Fallacy: Why Transformer Depth Limits Are Not Physical Laws. \emph{Unpublished manuscript}.
\bibitem[Hossenfelder(2026b)]{hossenfelder2026_holographic} Hossenfelder, S. (2026b). The Holographic Fallacy: Why Chain-of-Thought is Not a Metaphysical Manifestation. \emph{Unpublished manuscript}.
\end{thebibliography}

\end{document}
